\clearpage\thispagestyle{empty}
\pdfbookmark[0]{Quellen zu Gauß an Bode (10. August 1816)}{bode1816aug-sources}
\begin{center}\textbf{Quellen zu Gauß an Bode (10. August 1816)}\end{center}
\noindent Diese KI-Übertragung ins heutige Deutsch folgt dem Brief von Carl Friedrich Gauß an Johann Elert Bode in den \href{https://gdz.sub.uni-goettingen.de/id/PPN236007467}{\textit{Werken}, Band~VI} (Göttingen, 1874), S.~392--393. Das \textit{Astronomische Jahrbuch für 1819} (Berlin, 1816), S.~\href{https://api.digitale-sammlungen.de/iiif/image/v2/bsb10538348_00225/full/1600,/0/default.jpg}{219}--\href{https://api.digitale-sammlungen.de/iiif/image/v2/bsb10538348_00227/full/1600,/0/default.jpg}{221}, druckt den Brief unter einer Überschrift zur Opposition der Pallas für 1817, zu den Beobachtungen und neuen Bahnelementen der Juno und zu veränderlichen Sternen. Die Jahrbuchlinks führen zum Exemplar der Bayerischen Staatsbibliothek.
\par\smallskip\noindent Das Jahrbuch enthält zwei Fußnoten. Bei den beiliegenden Ephemeriden auf S.~\href{https://api.digitale-sammlungen.de/iiif/image/v2/bsb10538348_00225/full/1600,/0/default.jpg}{219} steht \textit{Folgen nachher.} („Sie folgen weiter unten.“). Bei Hardings Ephemeride auf S.~\href{https://api.digitale-sammlungen.de/iiif/image/v2/bsb10538348_00227/full/1600,/0/default.jpg}{221} steht \textit{Folgt nachher.} („Sie folgt weiter unten.“). Der Brief in der Werkausgabe endet mit der letzten jährlichen Abnahme in der Tabelle zum veränderlichen Stern. Die beigelegten Ephemeriden gehören dort nicht zum Brieftext.
\par\smallskip\noindent Auf S.~\href{https://api.digitale-sammlungen.de/iiif/image/v2/bsb10538348_00226/full/1600,/0/default.jpg}{220} fehlt in der Tabelle zur Lichtstärke der Eintrag für 1818. In der Werkausgabe steht 0.02122. Das Jahrbuch gibt in der Beobachtung vom 12.~Juni $27''{,}5$ an und für die tägliche tropische Bewegung $812''{,}40187$. In der Werkausgabe stehen $27^{\mathrm s}5$ und $812''40187$. KI-Übertragung und KI-Satz.
\clearpage
